% sections/introduction.tex
Colour fundus photographs show the retinal vessels directly, and changes in these vessels accompany common eye and systemic
diseases: diabetic retinopathy is a microvascular disease whose severe stage involves the growth of new retinal
vessels~\cite{wong2016dr}, and hypertension is associated with changes in the retinal vessels~\cite{wong2004hypertensive}.
Automated segmentation of the vessel tree is a building block for measuring such changes~\cite{abramoff2010retinal}, and
many methods have been proposed for it~\cite{fraz2012survey}.

The vessels at the ends of the tree are a few pixels wide and have low contrast. In the training masks used here,
centreline points with a local vessel radius below 3~pixels at 2048 $\times$ 2048 make up \ShareThin\% of the vessel
centreline. Such vessels contain few pixels, so pixel-overlap scores such as the Dice coefficient~\cite{dice1945}, which
weight every pixel equally, change little when many of them are missed. Validation guidelines therefore recommend metrics
that match the property of interest~\cite{maierhein2024metrics,reinke2024pitfalls}, such as clDice for tubular
structures~\cite{shit2021cldice}.

Two kinds of method are often proposed for fine vessels. One adds modules such as attention
gates~\cite{oktay2018attention} and multi-scale blocks~\cite{yu2016dilated,liu2022msunet,cao2023mfaunet} to the
U-Net~\cite{ronneberger2015unet}, a convolutional neural network (CNN); the other fine-tunes the Segment Anything Model
(SAM)~\cite{kirillov2023sam}, which fails on vessels when used zero-shot~\cite{shi2023samzeroshot,qiu2023learnable}. Whether either helps thin vessels is hard to judge from published
results, which come from different protocols and are often summarized with overall accuracy or
Dice~\cite{liu2022msunet,cao2023mfaunet}; variation between training runs alone can change the outcome of a
comparison~\cite{bouthillier2021variance}. On FIVES~\cite{jin2022fives}, earlier work compared
architectures~\cite{fadugba2024benchmark} and input resolutions~\cite{amarnath2026resolution}; the present study
combines thin-vessel metrics with a factorial design, several training seeds, an analysis of threshold choice and a
SAM configuration matched to the CNNs' output grid.

Resolution conditions the questions. At the 512 $\times$ 512 input of our CNNs, a thin vessel is at most about one pixel
wide, whereas SAM's decoder predicts a 256 $\times$ 256 grid, so each of its output pixels covers 8 $\times$ 8 pixels of the
original image. Within this setting we asked three questions. (1) Do attention gates and multi-scale blocks improve
thin-vessel segmentation when everything else is held fixed? (2) Does a fine-tuned SAM offer an advantage, and how do a
whole-image configuration and a tiled configuration with the CNNs' output grid differ? (3) Do overlap metrics reveal the
differences between these methods?

The study makes five contributions:
\begin{itemize}
\item An audit of FIVES, confirmed on the official release, which found duplicate images with different masks, empty masks, one eye in both
the training and test sets, annotations outside the eye and burned-in text, and a grouped split that accounts for them.
\item A 2 $\times$ 2 ablation of attention gates and multi-scale blocks on a fixed U-Net backbone, with three training seeds
per model and identical data, training and stopping rules.
\item Prompt-free fine-tuning of SAM on whole images and on full-resolution tiles whose combined output matches the CNNs'
grid, with a check of the tile seams.
\item Evaluation at full resolution with centreline recall by vessel width, clDice and normalized surface Dice, and
mask-resampling references for each output grid.
\item Separate reporting of training-seed and image-level variability, an analysis of threshold sensitivity, and analysis
decisions recorded in a public commit before the test set was used.
\end{itemize}
